\customchapter{INTRODUCCIÓN} 

\introsection{Presentación del tema de investigación}

Los modelos de aprendizaje automático ---o \textit{Machine Learning} (ML)--- desplegados
en producción operan bajo la suposición implícita de que la distribución de los datos de
entrada permanece estable a lo largo del tiempo. Sin embargo, en entornos reales esta
suposición rara vez se sostiene: los patrones demográficos evolucionan, los comportamientos
de los usuarios cambian y los contextos operativos se transforman de manera continua
\cite{mahadevan2024cost}. Cuando la distribución de los datos se aleja de la distribución de
entrenamiento, el rendimiento del modelo se degrada de forma silenciosa pero sistemática,
un fenómeno conocido como \textit{distribution shift} o desplazamiento de la distribución.
 
En \textit{pipelines} de ML de dos etapas ---compuestos por una etapa de \textit{feature
engineering} y una etapa de selección y entrenamiento del modelo predictivo--- esta
degradación puede originarse en cualquiera de las dos etapas o en ambas simultáneamente.
La etapa de preprocesamiento aprende estadísticas de la distribución de entrenamiento:
medias, desviaciones estándar, rangos de valores y frecuencias de categorías. Si la
distribución de los datos de entrada cambia (\textit{covariate shift}), estas estadísticas
quedan desactualizadas y el modelo recibe representaciones distorsionadas de los datos
reales. Por otro lado, si la relación entre las variables de entrada y la variable objetivo
cambia (\textit{concept drift}), el modelo predictivo habrá aprendido reglas que ya no
describen el mundo correctamente, independientemente de lo bien que el preprocesamiento
transforme los datos \cite{cai2023disde}.
 
La respuesta convencional ante la degradación del rendimiento es el reentrenamiento
completo del \textit{pipeline}, que resulta costoso en términos computacionales y
operativos \cite{mahadevan2024cost}. Esta práctica ignora una distinción fundamental: si el
origen del problema es únicamente la etapa de preprocesamiento, actualizar solo esa etapa
podría recuperar el rendimiento a una fracción del costo del reentrenamiento completo. De
manera análoga, si el problema radica en el modelo predictivo, reentrenar únicamente esa
etapa puede ser suficiente. La pregunta de cuándo cada estrategia de actualización selectiva
es suficiente, y a qué costo, no ha sido respondida de forma sistemática en la literatura
existente.
 
Resolver este problema requiere no solo diagnosticar el tipo y la severidad del drift
observado, sino también aprender una política de decisión que, ante un nuevo escenario de
degradación, seleccione la acción más eficiente. Esta investigación propone DARL
(\textit{Drift-Aware Reinforcement Learner}), un framework que formaliza esa política como
un agente de aprendizaje por refuerzo entrenado sobre un Proceso de Decisión de Markov
Parcialmente Observable (POMDP), donde el agente aprende a partir de la experiencia
acumulada sobre múltiples episodios de drift qué acción de actualización maximiza la
recuperación del rendimiento al menor costo computacional.


\introsection{Descripción de la situación problemática}

Los \textit{pipelines} de ML sobre datos tabulares se emplean de manera extensa en
dominios críticos como la salud, las finanzas y las políticas públicas \cite{gardner2023tableshift}.
En estos contextos, la degradación del rendimiento del modelo tiene consecuencias directas
sobre la calidad de las decisiones: un modelo de predicción de mortalidad hospitalaria que
se deteriora silenciosamente puede comprometer la asignación de recursos clínicos, y un
modelo de predicción de ingresos sesgado por cambios demográficos puede distorsionar
políticas de asistencia social.
 
La literatura reciente ha avanzado en el diagnóstico del tipo de drift que afecta a un
modelo \cite{cai2023disde,singh2025shift}, en el estudio de la robustez de distintos modelos frente a
\textit{distribution shift} en \textit{benchmarks} tabulares \cite{gardner2023tableshift} y en la
decisión de cuándo reentrenar considerando el costo computacional \cite{mahadevan2024cost,
shakhovska2025severity}. Sin embargo, estos trabajos presentan una limitación compartida: tratan el
modelo como una unidad indivisible y no estudian la posibilidad de actualizar selectivamente
solo una de las dos etapas del \textit{pipeline}.
 
Esta brecha tiene consecuencias prácticas importantes. En producción, los equipos de
\textit{Data Science} enfrentan la decisión de qué actualizar ante una caída de rendimiento
detectada, sin disponer de un mecanismo que aprenda de manera sistemática cuál es la acción
óptima en función del tipo y la severidad de la deriva observada. La ausencia de una
política de decisión aprendida ---que mapee el estado de las distribuciones observadas a la
acción de actualización más eficiente--- obliga a tomar decisiones \textit{ad hoc},
frecuentemente costosas e innecesarias.


\introsection{Formulación del problema}

La degradación del rendimiento de \textit{pipelines} de ML de dos etapas bajo
\textit{distribution shift} puede originarse en la etapa de \textit{feature engineering},
en la etapa del modelo predictivo o en ambas simultáneamente. Aunque existen métodos para
diagnosticar el tipo de drift \cite{cai2023disde,singh2025shift} y criterios para decidir cuándo
reentrenar \cite{mahadevan2024cost}, no existe un mecanismo que aprenda de forma sistemática
cuál de las cuatro posibles acciones de actualización ---no hacer nada, actualizar solo el
preprocesamiento, actualizar solo el modelo, o reentrenar el \textit{pipeline} completo---
produce la mejor recuperación del rendimiento al menor costo computacional, en función del
tipo de shift (\textit{covariate shift} o \textit{concept drift}) y su severidad (leve,
moderada o severa) sobre datos tabulares reales.
 
A partir de ello, se plantea la siguiente pregunta de investigación: ¿Es posible aprender
una política de decisión mediante aprendizaje por refuerzo que, a partir de señales de
monitoreo del estado de las distribuciones de un \textit{pipeline} de ML de dos etapas,
seleccione la estrategia de actualización selectiva que maximiza la recuperación del
rendimiento predictivo al menor costo computacional, generalizando a escenarios de drift
no vistos durante el entrenamiento?


\introsection{Objetivos de investigación}

Diseñar, implementar y evaluar empíricamente DARL (\textit{Drift-Aware
Reinforcement Learner}), un framework de actualización selectiva para \textit{pipelines} de
ML de dos etapas sobre datos tabulares que aprende una política de decisión mediante
aprendizaje por refuerzo sobre un POMDP formulado a partir del tipo y la severidad del
\textit{distribution shift} detectado, y que cuantifica la recuperación de rendimiento y el
costo computacional asociado a cada acción de actualización.
 
\section*{Objetivos específicos}
 
\begin{enumerate}
    \item Implementar un módulo de monitoreo que detecte y clasifique el tipo de
    \textit{data drift} ---\textit{covariate shift} y \textit{concept drift}---
    mediante el índice de estabilidad poblacional (PSI), el test de
    Kolmogorov--Smirnov (KS) y el \textit{Classifier Two-Sample Test} (C2ST),
    y que cuantifique su severidad mediante un \textit{severity score} combinado.
    Estas métricas constituyen el vector de observación del agente de aprendizaje
    por refuerzo.
 
    \item Formalizar el problema de decisión de actualización selectiva como un
    Proceso de Decisión de Markov Parcialmente Observable (POMDP), donde el espacio
    de acciones comprende las cuatro estrategias de actualización ---diferir la acción,
    actualizar el preprocesamiento, actualizar el modelo predictivo y reentrenar el
    \textit{pipeline} completo---, y la recompensa combina la recuperación del
    AUC-ROC y el costo computacional relativo de cada acción.
 
    \item Entrenar un agente de aprendizaje por refuerzo mediante el algoritmo
    \textit{Proximal Policy Optimization} (PPO) sobre episodios de drift sintético
    controlado, de modo que aprenda una política óptima $\pi^*$ que maximice la
    recompensa acumulada esperada sobre dos conjuntos de datos del
    \textit{benchmark} TableShift \cite{gardner2023tableshift}: Hospital Readmission
    (\texttt{diabetes\_readmission}) y PhysioNet/sepsis (\texttt{physionet}).
 
    \item Evaluar la política aprendida bajo condiciones de drift sintético de tipo
    y severidad conocidos, midiendo la recuperación del AUC-ROC y el costo
    computacional relativo de cada acción, y comparar su desempeño contra una
    tabla de decisión empírica construida como \textit{baseline}.
\end{enumerate}


\introsection{Justificación}

Desde una perspectiva científica, la literatura sobre \textit{distribution shift} en
aprendizaje automático ha avanzado de forma desigual. Los métodos de diagnóstico ---como
DISDE \cite{cai2023disde} y SHIFT \cite{singh2025shift}--- identifican con rigor el tipo de drift que
afecta a un modelo, pero no prescriben ninguna acción correctiva sobre el
\textit{pipeline}. Los métodos de decisión de reentrenamiento ---como CARA
\cite{mahadevan2024cost} y el framework de severidad de Shakhovska y Pukach
\cite{shakhovska2025severity}--- razonan sobre cuándo reentrenar considerando el costo, pero asumen
que el \textit{pipeline} es una unidad monolítica y no distinguen entre sus etapas. Los
trabajos sobre \textit{self-healing pipelines} \cite{tanna2025selfhealing} proponen estrategias de
remediación automática, pero no las evalúan empíricamente en función del tipo de deriva.
Esta investigación cierra ese ciclo al proponer DARL, un framework que conecta el
diagnóstico del drift con una política de decisión aprendida que selecciona la acción
correctiva y cuantifica sus consecuencias sobre el rendimiento y el costo, una contribución
que ninguno de los trabajos previos realiza de forma sistemática.
 
Desde una perspectiva práctica, los equipos de MLOps en producción frecuentemente
reentrenan el \textit{pipeline} completo ante cualquier degradación de rendimiento, sin
evaluar si una actualización parcial sería suficiente. Esta investigación proporciona un
agente entrenado que aprende a tomar decisiones más eficientes, reduciendo el tiempo de
cómputo y los costos operativos en escenarios donde el reentrenamiento selectivo es
suficiente para recuperar el rendimiento. Adicionalmente, el uso de aprendizaje por
refuerzo sobre un dominio de mantenimiento de modelos ---en contraste con su aplicación
habitual en juegos o robótica--- constituye una contribución metodológica con potencial de
generalización a otros problemas de gestión de \textit{pipelines} en producción.


\introsection{Alcance}

La investigación se circunscribe a \textit{pipelines} de ML de dos etapas 
(preprocesamiento y modelo predictivo) aplicados a tareas de clasificación binaria sobre
datos tabulares en modalidad \textit{batch}. La evaluación empírica se realiza sobre dos
conjuntos de datos del \textit{benchmark} TableShift: Hospital Readmission
(\texttt{diabetes\_readmission}) y PhysioNet/sepsis (\texttt{physionet})
\cite{gardner2023tableshift}. El \textit{distribution shift} se
inyecta de forma sintética y controlada en tres niveles de severidad (leve, moderada y
severa) para los tipos \textit{covariate shift}, \textit{concept drift} y su combinación.
El modelo predictivo empleado en los experimentos es XGBoost, con regresión logística como
línea base. El preprocesamiento sigue el esquema de la API de TableShift, extendido con
transformaciones de cuantiles para variables con distribución sesgada y codificación por
objetivo para variables categóricas de alta cardinalidad.
 
El agente de aprendizaje por refuerzo se plantea para entrenamiento mediante PPO,
con un entorno de simulación construido bajo una interfaz compatible con Gymnasium.
El entrenamiento se realiza exclusivamente sobre episodios
de drift sintético generados a partir de los dos datasets seleccionados; no se contempla
despliegue en entornos de producción en tiempo real ni integración con sistemas de
orquestación como Airflow o MLflow.


\introsection{Limitaciones y restricciones}

El framework DARL está orientado a uso académico y experimental. Los experimentos se ejecutan en
modalidad \textit{batch} y no cubren escenarios de \textit{streaming} con drift gradual y
continuo. La evaluación se restringe a los dos \textit{datasets} seleccionados y a un
modelo predictivo principal; la generalización de los resultados a otros dominios y
arquitecturas requiere validación adicional. Asimismo, se asume disponibilidad de etiquetas
en el conjunto de datos del dominio objetivo para realizar el reentrenamiento, condición que
en escenarios reales puede requerir anotación humana o mecanismos de \textit{delayed
feedback}. La política aprendida por el agente PPO está condicionada a la distribución de
episodios de entrenamiento generados; su transferencia a dominios con características de
drift muy distintas a las del entrenamiento requiere fine-tuning o reentrenamiento del
agente. Finalmente, el estudio no evalúa estrategias de defensa ante drift ni métodos de
reentrenamiento incremental, que se consideran líneas de trabajo futuro.
